% =============================================================================
%  Annexes
% =============================================================================

\chapter{Matrice des preuves}\label{annexe:preuves}

Chaque preuve est un fichier de \code{docs/preuves/} : sortie horodatée de commandes réelles
(\code{.txt}) ou capture (\code{.png}). La matrice ci-dessous est générée à partir du dossier
technique par \code{rapport/generer-annexes.py}.

{\small% Fichier généré par rapport/generer-annexes.py — ne pas modifier à la main
\begin{longtable}{@{}P{1.2cm}P{1.7cm}P{6.4cm}P{4.6cm}@{}}
\caption{Matrice des preuves (\code{docs/preuves/})}\label{tab:matrice}\\
\toprule \textbf{Preuve} & \textbf{Date} & \textbf{Contenu} & \textbf{Fichier(s)} \\ \midrule \endfirsthead
\toprule \textbf{Preuve} & \textbf{Date} & \textbf{Contenu} & \textbf{Fichier(s)} \\ \midrule \endhead
\bottomrule \endfoot
\preuve{P-01} & 19/10/2026 & VM Proxmox, snapshots, pont vmbr1, règles nftables & {\scriptsize \texttt{P-\allowbreak 01-\allowbreak proxmox-\allowbreak vms-\allowbreak et-\allowbreak reseau.txt}} \\
\preuve{P-02} & 19/10/2026 & LINUX01 : IP, connectivité, Internet, isolation + contre-test & {\scriptsize \texttt{P-\allowbreak 02-\allowbreak linux01-\allowbreak ip-\allowbreak connectivite-\allowbreak isolation.txt}} \\
\preuve{P-02b} & 19/10/2026 & OPS01 : IP, connectivité, Internet, isolation & {\scriptsize \texttt{P-\allowbreak 02b-\allowbreak ops01-\allowbreak ip-\allowbreak connectivite-\allowbreak isolation.txt}} \\
\preuve{P-03} & 19/10/2026 & Capture : WIN-DC01 installé, Server Manager & {\scriptsize \texttt{P-\allowbreak 03-\allowbreak win-\allowbreak dc01-\allowbreak installe-\allowbreak server-\allowbreak manager.png}} \\
\preuve{P-04} & 19/10/2026 & LINUX01 : DNS du domaine (A, PTR, SRV) & {\scriptsize \texttt{P-\allowbreak 04-\allowbreak linux01-\allowbreak resolution-\allowbreak dns-\allowbreak domaine.txt}} \\
\preuve{P-04b} & 19/10/2026 & OPS01 : DNS du domaine (A, PTR, SRV) & {\scriptsize \texttt{P-\allowbreak 04b-\allowbreak ops01-\allowbreak resolution-\allowbreak dns-\allowbreak domaine.txt}} \\
\preuve{P-04c} & 19/10/2026 & WIN-DC01 : IP, ping et résolution des serveurs & {\scriptsize \texttt{P-\allowbreak 04c-\allowbreak win-\allowbreak dc01-\allowbreak ip-\allowbreak dns.txt}} \\
\preuve{P-05} & 19/10/2026 & Arbre complet des snapshots Proxmox des 3 VM, snapshot final s05 & {\scriptsize \texttt{P-\allowbreak 05-\allowbreak proxmox-\allowbreak snapshots-\allowbreak arbre-\allowbreak complet.txt}} \\
\preuve{P-10a} & 19/10/2026 & 1er essai de promotion en échec (\code{reboot_timeout}) & {\scriptsize \texttt{P-\allowbreak 10a-\allowbreak adds-\allowbreak premier-\allowbreak essai-\allowbreak echec-\allowbreak reboot\_timeout.txt}} \\
\preuve{P-10} & 19/10/2026 & Installation AD DS + forêt infraready.test & {\scriptsize \texttt{P-\allowbreak 10-\allowbreak adds-\allowbreak installation-\allowbreak foret.txt}} \\
\preuve{P-11} & 19/10/2026 & OU, groupes, comptes AD & {\scriptsize \texttt{P-\allowbreak 11-\allowbreak ad-\allowbreak ou-\allowbreak groupes-\allowbreak comptes.txt}} \\
\preuve{P-12} & 19/10/2026 & État du domaine, comptes, DNS, zones, dcdiag & {\scriptsize \texttt{P-\allowbreak 12-\allowbreak domaine-\allowbreak dns-\allowbreak verifications.txt}} \\
\preuve{P-12b} & 19/10/2026 & dcdiag tests fonctionnels + NTP & {\scriptsize \texttt{P-\allowbreak 12b-\allowbreak dcdiag-\allowbreak tests-\allowbreak fonctionnels-\allowbreak ntp.txt}} \\
\preuve{P-13} & 19/10/2026 & Journaux Windows (Security, Directory Service, DNS, System) & {\scriptsize \texttt{P-\allowbreak 13-\allowbreak journaux-\allowbreak windows-\allowbreak ad.txt}} \\
\preuve{P-14} & 20/10/2026 & Windows : état initial, services, processus, journal System, script \code{Check-DnsClient.ps1} (0/1/2) & {\scriptsize \texttt{P-\allowbreak 14-\allowbreak win-\allowbreak dc01-\allowbreak processus-\allowbreak services-\allowbreak script-\allowbreak dnsclient.txt}} \\
\preuve{P-15} & 20/10/2026 & \code{windows_adds.yml} idempotent après correction du DNS du DC (\code{changed=0} ×2) & {\scriptsize \texttt{P-\allowbreak 15-\allowbreak windows-\allowbreak adds-\allowbreak dns-\allowbreak idempotence.txt}} \\
\preuve{P-20} & 19/10/2026 & Comptes, groupes, permissions Linux + tests de droits & {\scriptsize \texttt{P-\allowbreak 20-\allowbreak linux01-\allowbreak comptes-\allowbreak groupes-\allowbreak permissions.txt}} \\
\preuve{P-21} & 19/10/2026 & nginx, service et processus & {\scriptsize \texttt{P-\allowbreak 21-\allowbreak linux01-\allowbreak service-\allowbreak nginx-\allowbreak processus.txt}} \\
\preuve{P-22} & 19/10/2026 & Disque de test : contrôles, formatage, montage UUID & {\scriptsize \texttt{P-\allowbreak 22-\allowbreak linux01-\allowbreak disque-\allowbreak test-\allowbreak montage.txt}} \\
\preuve{P-23} & 19/10/2026 & Mise à jour contrôlée d'Ubuntu & {\scriptsize \texttt{P-\allowbreak 23-\allowbreak linux01-\allowbreak mise-\allowbreak a-\allowbreak jour-\allowbreak controlee.txt}} \\
\preuve{P-24} & 19/10/2026 & Analyse des journaux Linux & {\scriptsize \texttt{P-\allowbreak 24-\allowbreak linux01-\allowbreak analyse-\allowbreak journaux.txt}} \\
\preuve{P-25} & 20/10/2026 & Script \code{check-nginx.sh} testé dans 3 situations (codes 0/1/2) & {\scriptsize \texttt{P-\allowbreak 25-\allowbreak script-\allowbreak check-\allowbreak nginx-\allowbreak trois-\allowbreak situations.txt}} \\
\preuve{P-26} & 20/10/2026 & LINUX01 : état initial, processus, mémoire, disque, services, journaux nginx & {\scriptsize \texttt{P-\allowbreak 26-\allowbreak linux01-\allowbreak processus-\allowbreak services-\allowbreak journaux.txt}} \\
\preuve{P-30} & 19/10/2026 & Ansible : version, inventaire, connexions & {\scriptsize \texttt{P-\allowbreak 30-\allowbreak ansible-\allowbreak inventaire-\allowbreak connexions.txt}} \\
\preuve{P-31} & 19/10/2026 & Playbook Linux — 1er passage & {\scriptsize \texttt{P-\allowbreak 31-\allowbreak playbook-\allowbreak linux-\allowbreak 1er-\allowbreak passage.txt}} \\
\preuve{P-32} & 19/10/2026 & Playbook Linux — 2e passage (idempotence) + UFW & {\scriptsize \texttt{P-\allowbreak 32-\allowbreak playbook-\allowbreak linux-\allowbreak 2e-\allowbreak passage-\allowbreak idempotence.txt}} \\
\preuve{P-33} & 19/10/2026 & Playbook Windows — 1er passage & {\scriptsize \texttt{P-\allowbreak 33-\allowbreak playbook-\allowbreak windows-\allowbreak 1er-\allowbreak passage.txt}} \\
\preuve{P-34} & 19/10/2026 & Playbook Windows — 2e passage (idempotence) & {\scriptsize \texttt{P-\allowbreak 34-\allowbreak playbook-\allowbreak windows-\allowbreak 2e-\allowbreak passage-\allowbreak idempotence.txt}} \\
\preuve{P-35} & 19/10/2026 & Playbook Windows — correction PTR du DC & {\scriptsize \texttt{P-\allowbreak 35-\allowbreak playbook-\allowbreak windows-\allowbreak correction-\allowbreak ptr-\allowbreak dc.txt}} \\
\preuve{P-36} & 19/10/2026 & Diagnostic : enregistrements de la zone inverse (statique/dynamique) & {\scriptsize \texttt{P-\allowbreak 36-\allowbreak diagnostic-\allowbreak ptr-\allowbreak dc-\allowbreak statique-\allowbreak vs-\allowbreak dynamique.txt}} \\
\preuve{P-36b} & 19/10/2026 & Zone inverse après réenregistrement du DC & {\scriptsize \texttt{P-\allowbreak 36b-\allowbreak ptr-\allowbreak dc-\allowbreak apres-\allowbreak reenregistrement-\allowbreak dynamique.txt}} \\
\preuve{P-37} & 19/10/2026 & Playbook Windows stable après réenregistrement du DC (changed=0 ×2) & {\scriptsize \texttt{P-\allowbreak 37-\allowbreak playbook-\allowbreak windows-\allowbreak ptr-\allowbreak dc-\allowbreak stable.txt}} \\
\preuve{P-38} & 20/10/2026 & \code{linux_base.yml} + \code{site_message} et contrôle de contenu — 1ᵉʳ passage & {\scriptsize \texttt{P-\allowbreak 38-\allowbreak playbook-\allowbreak linux-\allowbreak site-\allowbreak message-\allowbreak 1er-\allowbreak passage.txt}} \\
\preuve{P-38b} & 20/10/2026 & \code{linux_base.yml} — 2ᵉ passage (\code{changed=0}) & {\scriptsize \texttt{P-\allowbreak 38b-\allowbreak playbook-\allowbreak linux-\allowbreak site-\allowbreak message-\allowbreak 2e-\allowbreak passage.txt}} \\
\preuve{P-39a} & 20/10/2026 & Simulation \code{--check --diff} en échec (blkid / uri non exécutés) & {\scriptsize \texttt{P-\allowbreak 39a-\allowbreak playbook-\allowbreak linux-\allowbreak mode-\allowbreak check-\allowbreak echec.txt}} \\
\preuve{P-39} & 20/10/2026 & Simulation corrigée (\code{check_mode: false}) puis exécution \code{changed=0} & {\scriptsize \texttt{P-\allowbreak 39-\allowbreak playbook-\allowbreak linux-\allowbreak mode-\allowbreak check-\allowbreak corrige.txt}} \\
\preuve{P-40} & 19/10/2026 & Maintenance Linux (linux01, ops01) : pré-contrôles, apt, post-contrôles & {\scriptsize \texttt{P-\allowbreak 40-\allowbreak linux-\allowbreak maintenance-\allowbreak standard.txt}} \\
\preuve{P-40b} & 19/10/2026 & Rapports de maintenance générés (\code{~/infraready-logs/}) & {\scriptsize \texttt{P-\allowbreak 40b-\allowbreak rapports-\allowbreak maintenance-\allowbreak generes.txt}} \\
\preuve{P-41a} & 19/10/2026 & Maintenance Windows — 1er essai : lecteur CD vu comme volume plein & {\scriptsize \texttt{P-\allowbreak 41a-\allowbreak windows-\allowbreak maintenance-\allowbreak 1er-\allowbreak essai-\allowbreak lecteur-\allowbreak cd.txt}} \\
\preuve{P-41} & 19/10/2026 & Maintenance Windows — recherche des mises à jour (5 trouvées) & {\scriptsize \texttt{P-\allowbreak 41-\allowbreak windows-\allowbreak maintenance-\allowbreak recherche-\allowbreak maj.txt}} \\
\preuve{P-42} & 19/10/2026 & Maintenance Windows — installation (5 MAJ, redémarrage du DC, boucle KB5007651) & {\scriptsize \texttt{P-\allowbreak 42-\allowbreak windows-\allowbreak maintenance-\allowbreak installation-\allowbreak maj.txt}} \\
\preuve{P-42b} & 19/10/2026 & Diagnostic WIN-DC01 : historique WU, journal Système, plateforme Defender & {\scriptsize \texttt{P-\allowbreak 42b-\allowbreak win-\allowbreak dc01-\allowbreak diagnostic-\allowbreak kb5007651-\allowbreak historique-\allowbreak wu.txt}} \\
\preuve{P-42c} & 19/10/2026 & Maintenance Windows — passage complet après exclusion de KB5007651 & {\scriptsize \texttt{P-\allowbreak 42c-\allowbreak windows-\allowbreak maintenance-\allowbreak installation-\allowbreak complete-\allowbreak apres-\allowbreak exclusion.txt}} \\
\preuve{P-43} & 19/10/2026 & unattended-upgrades : configuration, minuteries, simulation & {\scriptsize \texttt{P-\allowbreak 43-\allowbreak linux01-\allowbreak mises-\allowbreak a-\allowbreak jour-\allowbreak automatiques-\allowbreak unattended.txt}} \\
\preuve{P-70} & 19/10/2026 & Incident nginx : reproduction contrôlée et constat du symptôme (404) & {\scriptsize \texttt{P-\allowbreak 70-\allowbreak incident-\allowbreak nginx-\allowbreak reproduction-\allowbreak et-\allowbreak symptome.txt}} \\
\preuve{P-71} & 19/10/2026 & Incident nginx : diagnostic par élimination, cause confirmée & {\scriptsize \texttt{P-\allowbreak 71-\allowbreak incident-\allowbreak nginx-\allowbreak diagnostic.txt}} \\
\preuve{P-72a} & 19/10/2026 & Incident nginx : 1ᵉʳ essai de correction, handler non exécuté & {\scriptsize \texttt{P-\allowbreak 72a-\allowbreak incident-\allowbreak nginx-\allowbreak correction-\allowbreak 1er-\allowbreak essai-\allowbreak handler-\allowbreak non-\allowbreak execute.txt}} \\
\preuve{P-72} & 19/10/2026 & Incident nginx : correction avec le playbook corrigé (reproduction n°2) & {\scriptsize \texttt{P-\allowbreak 72-\allowbreak incident-\allowbreak nginx-\allowbreak correction-\allowbreak playbook-\allowbreak corrige.txt}} \\
\preuve{P-73} & 19/10/2026 & Incident nginx : retour à la normale et 2ᵉ passage \code{changed=0} & {\scriptsize \texttt{P-\allowbreak 73-\allowbreak incident-\allowbreak nginx-\allowbreak retour-\allowbreak normale-\allowbreak et-\allowbreak second-\allowbreak passage.txt}} \\
\preuve{P-74} & 20/10/2026 & Scénario 1 : variable incorrecte & {\scriptsize \texttt{P-\allowbreak 74-\allowbreak scenario1-\allowbreak variable-\allowbreak incorrecte.txt}} \\
\preuve{P-75} & 20/10/2026 & Scénario 2 : variable Jinja2 inexistante & {\scriptsize \texttt{P-\allowbreak 75-\allowbreak scenario2-\allowbreak variable-\allowbreak jinja2-\allowbreak inexistante.txt}} \\
\preuve{P-76} & 20/10/2026 & Scénario 3 : mauvais chemin de destination & {\scriptsize \texttt{P-\allowbreak 76-\allowbreak scenario3-\allowbreak mauvais-\allowbreak chemin-\allowbreak de-\allowbreak destination.txt}} \\
\preuve{P-77} & 20/10/2026 & Scénario 4 : handler non notifié & {\scriptsize \texttt{P-\allowbreak 77-\allowbreak scenario4-\allowbreak handler-\allowbreak non-\allowbreak notifie.txt}} \\
\preuve{P-78} & 20/10/2026 & Scénario 5 : nginx arrêté & {\scriptsize \texttt{P-\allowbreak 78-\allowbreak scenario5-\allowbreak nginx-\allowbreak arrete.txt}} \\
\preuve{P-79} & 20/10/2026 & Scénario 6 : contrôle fonctionnel du contenu & {\scriptsize \texttt{P-\allowbreak 79-\allowbreak scenario6-\allowbreak controle-\allowbreak fonctionnel-\allowbreak contenu.txt}} \\
\preuve{P-80} & 23/10/2026 & Séance 6 : état initial des 3 serveurs avant toute modification, blocages initiaux & {\scriptsize \texttt{P-\allowbreak 80-\allowbreak seance6-\allowbreak etat-\allowbreak initial-\allowbreak trois-\allowbreak serveurs.txt} \newline \texttt{P-\allowbreak 80-\allowbreak rendu-\allowbreak win-\allowbreak dc01-\allowbreak evenements-\allowbreak arret.png}} \\
\preuve{P-80c} & 23/10/2026 & Arbre des snapshots des 3 VM (API Proxmox), dont \code{s09_avant_seance6} & {\scriptsize \texttt{P-\allowbreak 80c-\allowbreak proxmox-\allowbreak snapshots-\allowbreak s09.txt} \newline \texttt{P-\allowbreak 80c-\allowbreak rendu-\allowbreak proxmox-\allowbreak snapshots-\allowbreak s09.png}} \\
\preuve{P-80b} & 23/10/2026 & \code{openipmi} désactivé par Ansible ; contrôle de santé sans ECHEC & {\scriptsize \texttt{P-\allowbreak 80b-\allowbreak seance6-\allowbreak correction-\allowbreak openipmi-\allowbreak et-\allowbreak controle-\allowbreak sante-\allowbreak vert.txt}} \\
\preuve{P-81a} & 23/10/2026 & Redéploiement de la supervision : écart de règles corrigé, vérification des cibles instable & {\scriptsize \texttt{P-\allowbreak 81a-\allowbreak seance6-\allowbreak redeploiement-\allowbreak monitoring-\allowbreak verification-\allowbreak cibles-\allowbreak instable.txt}} \\
\preuve{P-81b} & 23/10/2026 & Vérification des cibles corrigée (attente du 1ᵉʳ scrape), idempotence & {\scriptsize \texttt{P-\allowbreak 81b-\allowbreak seance6-\allowbreak monitoring-\allowbreak verification-\allowbreak cibles-\allowbreak corrigee.txt}} \\
\preuve{P-81} & 23/10/2026 & Cibles attendues / observées, endpoints, métriques Linux et Windows, règles & {\scriptsize \texttt{P-\allowbreak 81-\allowbreak seance6-\allowbreak collecte-\allowbreak metriques-\allowbreak linux-\allowbreak windows.txt} \newline \texttt{P-\allowbreak 81-\allowbreak prometheus-\allowbreak cibles-\allowbreak up.png}} \\
\preuve{P-82} & 23/10/2026 & \code{controle_sante.yml} : exécution, \code{--check}, hôte injoignable, correctif \code{--limit} & {\scriptsize \texttt{P-\allowbreak 82-\allowbreak controle-\allowbreak sante-\allowbreak ansible-\allowbreak lecture-\allowbreak seule.txt} \newline \texttt{P-\allowbreak 82-\allowbreak rendu-\allowbreak tableau-\allowbreak controle-\allowbreak sante.png}} \\
\preuve{P-83} & 23/10/2026 & Tâche planifiée : timer systemd, journal, métriques textfile, 2 nouvelles règles & {\scriptsize \texttt{P-\allowbreak 83-\allowbreak controle-\allowbreak planifie-\allowbreak timer-\allowbreak systemd.txt} \newline \texttt{P-\allowbreak 83-\allowbreak rendu-\allowbreak timer-\allowbreak journal-\allowbreak metriques.png} \newline \texttt{P-\allowbreak 83-\allowbreak prometheus-\allowbreak regles-\allowbreak chargees.png}} \\
\preuve{P-84} & 23/10/2026 & Alertes \code{SiteWebIndisponible} et \code{ControlePlanifieEnRetard} : déclenchement et retour à la normale + capture & {\scriptsize \texttt{P-\allowbreak 84-\allowbreak alerte-\allowbreak sitewebindisponible-\allowbreak test-\allowbreak controle.txt} \newline \texttt{P-\allowbreak 84-\allowbreak prometheus-\allowbreak alerte-\allowbreak sitewebindisponible-\allowbreak firing.png} \newline \texttt{P-\allowbreak 84-\allowbreak grafana-\allowbreak dashboard-\allowbreak site-\allowbreak en-\allowbreak echec.png} \newline \texttt{P-\allowbreak 84-\allowbreak grafana-\allowbreak dashboard-\allowbreak controle-\allowbreak en-\allowbreak retard.png} \newline \texttt{P-\allowbreak 84-\allowbreak prometheus-\allowbreak historique-\allowbreak controle-\allowbreak succes.png} \newline \texttt{P-\allowbreak 84-\allowbreak prometheus-\allowbreak historique-\allowbreak alertes.png}} \\
\preuve{P-85} & 23/10/2026 & Dashboard Grafana v3 : déploiement, vérification par l'API + capture & {\scriptsize \texttt{P-\allowbreak 85-\allowbreak grafana-\allowbreak dashboard-\allowbreak v3-\allowbreak deploiement.txt} \newline \texttt{P-\allowbreak 85-\allowbreak grafana-\allowbreak dashboard-\allowbreak v3-\allowbreak exploitation.png}} \\
\preuve{P-86} & 23/10/2026 & Incident INC-2026-10-23-01 : arrêt non propre de WIN-DC01, diagnostic et reproduction & {\scriptsize \texttt{P-\allowbreak 86-\allowbreak incident-\allowbreak arret-\allowbreak non-\allowbreak propre-\allowbreak win-\allowbreak dc01-\allowbreak diagnostic.txt} \newline \texttt{P-\allowbreak 86-\allowbreak rendu-\allowbreak proxmox-\allowbreak arrets-\allowbreak vm100-\allowbreak reproduction.png}} \\
\preuve{P-87} & 23/10/2026 & Maintenance LINUX01 avec contrôle de santé avant / après & {\scriptsize \texttt{P-\allowbreak 87-\allowbreak maintenance-\allowbreak linux01-\allowbreak controle-\allowbreak avant-\allowbreak apres.txt} \newline \texttt{P-\allowbreak 87-\allowbreak rendu-\allowbreak maintenance-\allowbreak avant-\allowbreak apres.png}} \\
\preuve{P-88} & 23/10/2026 & Espace disque : observation ponctuelle et historique Prometheus & {\scriptsize \texttt{P-\allowbreak 88-\allowbreak observation-\allowbreak ponctuelle-\allowbreak vs-\allowbreak historique-\allowbreak disque.txt} \newline \texttt{P-\allowbreak 88-\allowbreak grafana-\allowbreak evolution-\allowbreak stockage-\allowbreak 19-\allowbreak 23-\allowbreak oct.png}} \\
\preuve{P-89a} & 23/10/2026 & Quotas — 1ᵉʳ essai non idempotent (\code{repquota} sans \code{-v}) & {\scriptsize \texttt{P-\allowbreak 89a-\allowbreak quotas-\allowbreak 1er-\allowbreak essai-\allowbreak non-\allowbreak idempotent.txt}} \\
\preuve{P-89} & 23/10/2026 & Quotas sur \code{/data} : idempotence, dépassement souple / strict, nettoyage & {\scriptsize \texttt{P-\allowbreak 89-\allowbreak quotas-\allowbreak disque-\allowbreak test-\allowbreak data-\allowbreak linux01.txt} \newline \texttt{P-\allowbreak 89-\allowbreak rendu-\allowbreak test-\allowbreak quotas-\allowbreak alice.png}} \\
\preuve{P-50a} & 19/10/2026 & Monitoring — 1er essai : échec de connexion admin Grafana & {\scriptsize \texttt{P-\allowbreak 50a-\allowbreak ops01-\allowbreak grafana-\allowbreak 1er-\allowbreak essai-\allowbreak echec-\allowbreak mot-\allowbreak de-\allowbreak passe-\allowbreak ini.txt}} \\
\preuve{P-50} & 19/10/2026 & Playbook Prometheus + Grafana (cibles up, admin OK) & {\scriptsize \texttt{P-\allowbreak 50-\allowbreak ops01-\allowbreak prometheus-\allowbreak grafana-\allowbreak installation.txt}} \\
\preuve{P-50b} & 19/10/2026 & Playbook monitoring — 2ᵉ passage (changed=0) & {\scriptsize \texttt{P-\allowbreak 50b-\allowbreak ops01-\allowbreak monitoring-\allowbreak 2e-\allowbreak passage-\allowbreak idempotence.txt}} \\
\preuve{P-51} & 19/10/2026 & Capture : cibles Prometheus (4/4 up) & {\scriptsize \texttt{P-\allowbreak 51-\allowbreak prometheus-\allowbreak targets.png}} \\
\preuve{P-52} & 19/10/2026 & Cibles, règles, métriques réelles, exporters, ports & {\scriptsize \texttt{P-\allowbreak 52-\allowbreak prometheus-\allowbreak cibles-\allowbreak metriques-\allowbreak regles.txt}} \\
\preuve{P-53} & 19/10/2026 & Capture : dashboard Grafana, état normal & {\scriptsize \texttt{P-\allowbreak 53-\allowbreak grafana-\allowbreak dashboard-\allowbreak infraready.png}} \\
\preuve{P-54} & 19/10/2026 & Démonstration InstanceDown : chronologie + capture Prometheus & {\scriptsize \texttt{P-\allowbreak 54-\allowbreak alerte-\allowbreak instancedown-\allowbreak demonstration.txt} \newline \texttt{P-\allowbreak 54-\allowbreak prometheus-\allowbreak alerte-\allowbreak instancedown.png}} \\
\preuve{P-55} & 19/10/2026 & Capture : dashboard pendant l'alerte (LINUX01 hors ligne) & {\scriptsize \texttt{P-\allowbreak 55-\allowbreak grafana-\allowbreak dashboard-\allowbreak alerte-\allowbreak linux01-\allowbreak hors-\allowbreak ligne.png}} \\
\preuve{P-56} & 20/10/2026 & Dashboard Grafana v2 (variable Hôte, Uptime) déployé, idempotence, requêtes vérifiées & {\scriptsize \texttt{P-\allowbreak 56-\allowbreak grafana-\allowbreak dashboard-\allowbreak v2-\allowbreak deploiement-\allowbreak idempotence.txt}} \\
\preuve{P-60a} & 19/10/2026 & GitHub Actions sur le dépôt de travail initial : 1er run en échec, 2ᵉ vert & {\scriptsize \texttt{P-\allowbreak 60a-\allowbreak github-\allowbreak actions-\allowbreak depot-\allowbreak de-\allowbreak travail-\allowbreak initial.txt}} \\
\preuve{P-60} & 20/10/2026 & GitHub Actions sur \code{jolanallen/InfraReady} : run vert, jobs, déclencheurs, auteurs & {\scriptsize \texttt{P-\allowbreak 60-\allowbreak github-\allowbreak actions-\allowbreak pipeline-\allowbreak vert.txt}} \\
\preuve{P-61} & 20/10/2026 & Validation des PR \#2 et \#4 : \code{verif-idempotence.sh} sur \code{windows_base} et \code{linux_base} (IDEMPOTENT) & {\scriptsize \texttt{P-\allowbreak 61-\allowbreak validation-\allowbreak pr-\allowbreak idempotence-\allowbreak playbooks.txt}} \\
\preuve{P-62} & 20/10/2026 & Tests négatifs (non idempotent → 1, échec → 2, vault absent ou invalide) ; diagnostic du DNS de \code{windows_adds} & {\scriptsize \texttt{P-\allowbreak 62-\allowbreak validation-\allowbreak pr-\allowbreak tests-\allowbreak negatifs.txt}} \\
\preuve{P-63} & 20/10/2026 & Validation de la PR \#1 : calcul d'occupation disque de \code{linux_maintenance} (simulation) & {\scriptsize \texttt{P-\allowbreak 63-\allowbreak validation-\allowbreak pr1-\allowbreak maintenance-\allowbreak linux-\allowbreak calcul-\allowbreak disque.txt}} \\
\preuve{P-64} & 23/10/2026 & GitHub Actions sur les PR de la séance 6 : contrôles exécutés sur chaque branche ; PR \#8 à \#14 fusionnées & {\scriptsize \texttt{P-\allowbreak 64-\allowbreak github-\allowbreak actions-\allowbreak pr-\allowbreak seance6.txt} \newline \texttt{P-\allowbreak 64-\allowbreak github-\allowbreak pr-\allowbreak seance6-\allowbreak fusionnees.png} \newline \texttt{P-\allowbreak 64-\allowbreak github-\allowbreak actions-\allowbreak runs-\allowbreak verts.png}} \\
\end{longtable}
}

\chapter{Arborescence du dépôt}\label{annexe:arborescence}

\begin{footnotesize}
\begin{verbatim}
InfraReady/
├── .github/workflows/validation.yml     pipeline de validation (4 jobs)
├── ansible/
│   ├── ansible.cfg · requirements.yml
│   ├── inventory/hosts.yml
│   ├── inventory/group_vars/            all, linux, windows, monitoring (vaults chiffrés)
│   └── playbooks/                       windows_adds, windows_base, windows_maintenance,
│       │                                linux_base, linux_maintenance, linux_quotas,
│       │                                ops01_monitoring, controle_sante
│       └── templates/                   index.html.j2, nginx-infraready.conf.j2
├── lab/                                 modèles autoinstall / autounattend, réseau Proxmox
├── monitoring/
│   ├── prometheus/                      prometheus.yml, rules/infraready.rules.yml
│   └── grafana/                         dashboard infraready.json, provisioning
├── scripts/                             amorçage, vérification, diagnostic (Bash, PowerShell)
├── docs/
│   ├── procedure-maintenance.md · rapport-incident.md
│   └── preuves/                         P-01 … P-89 (texte horodaté et captures)
├── rapport/                             sources LaTeX de ce rapport, schémas, script de build
├── InfraReady - Dossier technique.md    dossier technique (Obsidian)
├── InfraReady - Séance 6.md             note de synthèse illustrée de la séance 6
├── InfraReady-Rapport-final.tex / .pdf  ce rapport
├── CLAUDE.md                            règles de contribution (Git, documentation, validation)
└── README.md
\end{verbatim}
\end{footnotesize}

\chapter{Journal des changements}\label{annexe:journal}

Changements appliqués au laboratoire, dans l'ordre chronologique, avec leur validation et le
moyen de retour arrière (\code{sNN_…} = snapshot Proxmox).

{\small% Fichier généré par rapport/generer-annexes.py — ne pas modifier à la main
\begin{longtable}{@{}P{2cm}P{6.3cm}P{3.4cm}P{2.6cm}@{}}
\caption{Journal des changements sur le laboratoire}\label{tab:journal}\\
\toprule \textbf{Date} & \textbf{Modification} & \textbf{Validation} & \textbf{Preuve / retour arrière} \\ \midrule \endfirsthead
\toprule \textbf{Date} & \textbf{Modification} & \textbf{Validation} & \textbf{Preuve / retour arrière} \\ \midrule \endhead
\bottomrule \endfoot
19/10/2026 14:42 & Pont \code{vmbr1} + NAT/isolation nftables sur l'hôte & \code{ifreload -a}, ping/isolation OK & {\small P-01, P-02 / \code{interfaces.bak-infraready}} \\
19/10/2026 14:47 & Installations automatiques des 3 VM & 3 VM joignables & {\small P-01, P-03 / —} \\
19/10/2026 14:58 & Snapshot de base & \code{qm listsnapshot} & {\small \code{s01_installation_ok}} \\
19/10/2026 15:05 & AD DS + forêt infraready.test & dcdiag fonctionnel OK & {\small P-10, P-12b / \code{s02_avant_adds}} \\
19/10/2026 15:14 & OU, groupes, comptes AD & Membres conformes & {\small P-11 / \code{s02_avant_adds}} \\
19/10/2026 15:20 & NTP externe sur le DC & Stratum 2 & {\small P-12b / —} \\
19/10/2026 15:23 & Comptes Linux, nginx, disque \code{/data} & Tests de droits OK & {\small P-20 à P-22 / \code{s02_avant_maj_c}} \\
19/10/2026 15:24 & \code{apt upgrade} LINUX01 (20 paquets) & Code retour 0 & {\small P-23 / \code{s02_avant_maj_c}} \\
19/10/2026 15:28 & Playbooks Linux et Windows (2 passages) & changed=0 au 2ᵉ passage & {\small P-31 à P-34 / \code{s03_avant_ansible}} \\
19/10/2026 & DNS des VM Linux → DC, PTR du DC & A/PTR/SRV résolus & {\small P-04, P-35 / \code{/root/50-cloud-init.yaml.bak-dns}} \\
19/10/2026 16:30 & Prometheus, exporters, Grafana sur OPS01 & Cibles UP, dashboard, alerte testée & {\small P-50 à P-55 / \code{s04_avant_maintenance_e}} \\
19/10/2026 16:48 & Maintenance Linux (LINUX01, OPS01), \code{unattended-upgrades} & Post-contrôles OK, reboot non requis & {\small P-40, P-43 / \code{s04_avant_maintenance_e}} \\
19/10/2026 16:54 & Windows Update sur WIN-DC01 (5 MAJ, redémarrage 18:04) & \code{failed=1} (boucle KB5007651) puis passage complet OK & {\small P-42 à P-42c / \code{s04_avant_maj_windows}} \\
19/10/2026 19:21 & \code{reject_list: [KB5007651]} + gabarit du rapport Windows & \code{failed=0}, rapport lisible & {\small P-42c / —} \\
19/10/2026 19:27 & Incident contrôlé : \code{root} nginx modifié sur LINUX01 & 404 constaté depuis OPS01 & {\small P-70 / \code{/root/infraready.nginx.bak-incident-h}} \\
19/10/2026 19:30 & \code{flush_handlers} avant les vérifications dans \code{linux_base.yml} ; rejeu sur LINUX01 & HTTP 200, \code{changed=0} au 2ᵉ passage & {\small P-72, P-73 / —} \\
19/10/2026 19:37 & Snapshot final des 3 VM & \code{qm listsnapshot} & {\small P-05 / \code{s05_phases_a_h_terminees}} \\
20/10/2026 12:50 & Snapshot des 3 VM avant la séance 4 & \code{qm snapshot} & {\small \code{s07_avant_seance4}} \\
20/10/2026 12:51 & \code{site_message} + contrôle HTTP et contenu dans \code{linux_base.yml} (2 passages) & \code{changed=1} puis \code{changed=0} & {\small P-38, P-38b / \code{s07_avant_seance4}} \\
20/10/2026 12:52 & Scripts \code{check-nginx.sh} et \code{Check-DnsClient.ps1} (3 situations chacun) & Codes 0 / 1 / 2 & {\small P-25, P-14 / —} \\
20/10/2026 12:57 & \code{check_mode: false} sur 4 tâches en lecture seule de \code{linux_base.yml} & Simulation \code{changed=0 failed=0} & {\small P-39a, P-39 / —} \\
20/10/2026 12:58 & 6 scénarios de panne contrôlée (7.4), un à la fois, jusqu'à 13:08 & Retour à l'état initial, \code{changed=0} & {\small P-74 à P-79 / \code{.etat-initial}, \code{s07_avant_seance4}} \\
20/10/2026 14:15 & Snapshot des 3 VM avant validation des PR \#1 à \#5 & \code{qm snapshot} & {\small \code{s08_avant_validation_pr3_5}} \\
20/10/2026 14:16 & Dashboard Grafana v2 déployé (\code{ops01_monitoring.yml}, 2 passages) & IDEMPOTENT, version 2 chargée & {\small P-56 / \code{s08_avant_validation_pr3_5}} \\
20/10/2026 14:18 & \code{windows_base.yml} et \code{linux_base.yml} via \code{verif-idempotence.sh} & IDEMPOTENT ×2 & {\small P-61 / —} \\
20/10/2026 14:21 & Tests négatifs ; exécution réelle non voulue de \code{windows_adds.yml} (DNS du DC rejoué, valeur inchangée) & Résolution DNS OK & {\small P-62 / \code{s08_avant_validation_pr3_5}} \\
20/10/2026 14:24 & \code{linux_maintenance.yml} en simulation (PR \#1) & Calcul identique à P-40 & {\small P-63 / —} \\
20/10/2026 14:26 & Correctif DNS de \code{windows_adds.yml} testé (PR \#6), puis copie d'OPS01 remise à la version de \code{main} & IDEMPOTENT & {\small P-15 (PR \#6) / \code{s08_avant_validation_pr3_5}} \\
23/10/2026 09:05 & Snapshot des 3 VM avant la séance 6 (API Proxmox) & Tâches \code{qmsnapshot} OK & {\small \code{s09_avant_seance6}} \\
23/10/2026 09:06 & État initial des 3 serveurs (lecture seule) & 5 écarts relevés & {\small P-80 / —} \\
23/10/2026 09:17 & Copie du projet sur OPS01 resynchronisée ; \code{ops01_monitoring.yml} rejoué (5 règles) puis corrigé (attente du 1ᵉʳ scrape) & IDEMPOTENT & {\small P-81a, P-81b / \code{s09_avant_seance6}} \\
23/10/2026 09:22 & Timer \code{infraready-controle-site} (\code{linux_base.yml}) et 2 règles d'alerte (\code{ops01_monitoring.yml}) & \code{changed=0} au 2ᵉ passage, 7 règles & {\small P-83 / \code{s09_avant_seance6}} \\
23/10/2026 09:26 & Test d'alerte : nginx arrêté sur LINUX01 puis redémarré (09:30) & firing à +3 min 08 s, résolue 27 s après & {\small P-84 / —} \\
23/10/2026 09:35 & \code{openipmi} désactivé (\code{linux_base.yml}, \code{ops01_monitoring.yml}) & Contrôle de santé sans ECHEC & {\small P-80b / \code{s09_avant_seance6}} \\
23/10/2026 09:38 & Reproduction : arrêt ACPI de WIN-DC01 demandé sans arrêt forcé & Ignoré pendant 90 s, VM restée allumée & {\small P-86 / —} \\
23/10/2026 09:43 & Dashboard Grafana v3 (\code{ops01_monitoring.yml}) & IDEMPOTENT, version 3 chargée & {\small P-85 / \code{s09_avant_seance6}} \\
23/10/2026 09:45 & Maintenance LINUX01 (3 paquets) entre deux contrôles de santé & Aucun écart de statut & {\small P-87 / \code{s09_avant_seance6}} \\
23/10/2026 09:48 & \code{/data} remonté avec \code{usrquota,grpquota} ; quotas alice / bruno (\code{linux_quotas.yml}) & IDEMPOTENT après correction ; dépassement testé & {\small P-89a, P-89 / \code{s09_avant_seance6}} \\
23/10/2026 09:54 & Test d'alerte : timer du contrôle planifié arrêté (09:54) puis redémarré (10:04) & \code{ControlePlanifieEnRetard} firing à 10:04:01, résolue à 10:04:17 & {\small P-84 / —} \\
\end{longtable}
}

\chapter{Glossaire}\label{annexe:glossaire}

\begin{description}[style=nextline,leftmargin=1.2cm]
  \item[AD DS] Active Directory Domain Services : annuaire Windows (comptes, groupes, ordinateurs, authentification Kerberos).
  \item[ACPI] Interface standard de gestion de l'alimentation ; l'hyperviseur l'utilise pour demander l'arrêt d'une VM.
  \item[Ansible Vault] Chiffrement AES256 de fichiers de variables Ansible contenant des secrets.
  \item[CI] Intégration continue : contrôles automatiques exécutés à chaque modification du dépôt (ici GitHub Actions).
  \item[DC] Contrôleur de domaine (WIN-DC01).
  \item[Exporter] Programme qui expose les métriques d'un système au format Prometheus (node\_exporter, windows\_exporter).
  \item[Gestionnaire (handler)] Tâche Ansible déclenchée seulement si une autre tâche a changé quelque chose (ex. recharger nginx).
  \item[Idempotence] Propriété d'une automatisation dont un second passage ne modifie plus rien (\code{changed=0}).
  \item[nftables] Pare-feu et traduction d'adresses du noyau Linux, successeur d'iptables.
  \item[Preuve P-xx] Fichier horodaté de \code{docs/preuves/} qui justifie un résultat de ce rapport.
  \item[PromQL] Langage de requête de Prometheus.
  \item[PTR] Enregistrement DNS inverse (adresse IP vers nom).
  \item[Scrape] Collecte périodique des métriques d'une cible par Prometheus (toutes les 15 s ici).
  \item[Snapshot] Instantané d'une VM permettant de revenir à un état antérieur.
  \item[Textfile (collecteur)] Fonction de node\_exporter qui publie des métriques écrites dans des fichiers \code{.prom}.
  \item[Timer systemd] Planification d'un service par systemd, équivalent moderne de cron.
  \item[UFW] Pare-feu simplifié d'Ubuntu.
  \item[WinRM] Protocole d'administration à distance de Windows, utilisé par Ansible.
\end{description}

\chapter{Références}\label{annexe:references}

\begin{itemize}
  \item Sujet du projet fil rouge InfraReady et supports pédagogiques des séances 4 (\emph{Stockage et
        processus}) et 6 (\emph{Maintenance opérationnelle et supervision}), Bachelor 3 — Administration
        systèmes Linux / Windows.
  \item Dossier technique du projet : \code{InfraReady - Dossier technique.md} ; procédure de
        maintenance : \code{docs/procedure-maintenance.md} ; rapports d'incident :
        \code{docs/rapport-incident.md}.
  \item Documentations officielles : Proxmox VE, Ansible (collections \code{ansible.windows},
        \code{microsoft.ad}, \code{community.general}), Prometheus, Grafana, node\_exporter,
        windows\_exporter, Microsoft Learn (AD DS, DNS, Windows Update), Ubuntu Server
        (autoinstall, unattended-upgrades, quota).
  \item Dépôt de remise : \url{https://github.com/jolanallen/InfraReady}.
\end{itemize}
